\section{Additional Generalization Experiments for Outcome Supervision}
\label{app:os_generalization}

We provide additional experiments to examine whether the optimization barrier of outcome supervision is specific to the backbone scale or the arithmetic reasoning setting used in the main experiments.
We first scale the backbone from GPT-2 to LLaMA-3.2-1B-Instruct on GSM8K-Aug.
We then evaluate the same supervision paradigms on ProntoQA~\citep{PrOntoQA}, a propositional logic QA benchmark requiring multi-hop deductive reasoning.
Together, these results show that the failure pattern of outcome-supervised Latent CoT is not confined to a small backbone or to arithmetic reasoning.

\subsection{Scale-Up Experiment on LLaMA-3.2-1B-Instruct}
\label{app:scale_up_os}

To examine whether the optimization barrier of outcome supervision is specific to the small GPT-2 backbone used in our main experiments, we conduct an additional scale-up experiment with LLaMA-3.2-1B-Instruct on GSM8K-Aug.
This setting uses a substantially larger and more modern backbone while keeping the same comparison between explicit CoT, outcome-supervised Latent CoT, and answer-only training.

We train all models with the AdamW optimizer using a learning rate of $1\mathrm{e}{-4}$, weight decay $0.01$, a cosine learning-rate schedule with warmup, gradient clipping at $1.0$, effective batch size $128$, and BF16 precision.
Experiments are run on $4$ RTX 5880 Ada GPUs.

\begin{table}[h]
\centering
\caption{Scale-up experiment on LLaMA-3.2-1B-Instruct with GSM8K-Aug. Outcome-supervised Latent CoT does not materially outperform the answer-only baseline, while explicit CoT remains substantially stronger.}
\label{tab:llama_scale_os}
\begin{tabular}{lcc}
\toprule
\textbf{Method} & \textbf{Best Acc.} & \textbf{Best Epoch} \\
\midrule
Explicit CoT & $60.20\%$ & $10$ \\
\textsc{OS-Latent} & $33.81\%$ & $9$ \\
\textsc{OS-No-CoT} & $32.98\%$ & $8$ \\
\bottomrule
\end{tabular}
\end{table}

As shown in Table~\ref{tab:llama_scale_os}, \textsc{OS-Latent} achieves $33.81\%$ best accuracy, which is close to the answer-only baseline at $32.98\%$ and far below explicit CoT at $60.20\%$.
This result reproduces the same qualitative pattern observed in the main GPT-2 experiments: adding latent states under outcome supervision alone does not reliably induce useful latent reasoning dynamics.
The substantial gap between explicit CoT and \textsc{OS-Latent} further indicates that the failure is not simply due to insufficient backbone capacity, since the same model benefits strongly from explicit intermediate reasoning supervision.

\subsection{Cross-Task Experiment on ProntoQA}
\label{app:prontoqa_os}

To examine whether the observed outcome-supervision failure is specific to arithmetic reasoning, we further evaluate on ProntoQA, a propositional logic QA benchmark requiring multi-hop deductive reasoning.
ProntoQA is structurally different from GSM8K-Aug: instead of numerical calculation, it requires the model to apply symbolic rules over a chain of logical implications.
This makes it a useful complementary testbed for evaluating whether latent CoT supervision patterns generalize beyond arithmetic.

\begin{table}[h]
\centering
\caption{Cross-task evaluation on ProntoQA and GSM8K-Aug. On ProntoQA, \textsc{OS-Latent} remains close to the answer-only baseline, while \textsc{PS-Latent} substantially improves performance. GSM8K-Aug results are included for comparison with the main experiments.}
\label{tab:prontoqa_generalization}
\begin{tabular}{lcc}
\toprule
\textbf{Method} & \textbf{ProntoQA Acc.} & \textbf{GSM8K-Aug Acc.} \\
\midrule
\textsc{PS-Latent} & $97.40\%$ & $31.2\%$ \\
Explicit CoT & $80.80\%$ & $43.1\%$ \\
\textsc{OS-Latent} & $78.60\%$ & $18.3\%$ \\
\textsc{OS-No-CoT} & $78.60\%$ & $18.7\%$ \\
\bottomrule
\end{tabular}
\end{table}

As shown in Table~\ref{tab:prontoqa_generalization}, \textsc{OS-Latent} obtains the same accuracy as the answer-only baseline on ProntoQA, indicating that latent states induced by outcome supervision alone do not provide a measurable benefit over direct answer prediction in this setting.
In contrast, \textsc{PS-Latent} reaches $97.40\%$ accuracy, substantially outperforming both \textsc{OS-Latent} and explicit CoT.
This result supports the same qualitative conclusion as the GSM8K-Aug experiments: outcome supervision alone is insufficient to reliably induce useful latent reasoning dynamics, whereas process supervision provides effective intermediate guidance.

We further evaluate the Unified Latent Probe (ULP) on ProntoQA.
The probe loss preserves the same information hierarchy observed in the main experiments: \textsc{PS-Latent} achieves lower probe loss than \textsc{OS-Latent} ($0.434$ vs. $0.546$).
Since lower ULP loss indicates better recovery of explicit reasoning information from latent states, this result provides additional evidence that the advantage of process supervision is reflected not only in task accuracy but also in the amount of recoverable reasoning information encoded in the latent trajectory.